\documentclass[11pt]{article}
\usepackage[margin=0.9in]{geometry}
\usepackage{amsmath,amssymb,amsthm,mathtools,microtype}
\usepackage{graphicx,xcolor}
\usepackage[colorlinks=true,linkcolor=blue!55!black,citecolor=blue!55!black,urlcolor=blue!55!black]{hyperref}

\newtheorem{theorem}{Theorem}
\newtheorem{corollary}[theorem]{Corollary}
\newtheorem{lemma}[theorem]{Lemma}
\newcommand{\Kop}{\mathcal K}

\title{Compact-Operator Algebras Determine the Underlying Banach Space}
\author{Candidate full solution to an explicit subquestion in arXiv:2505.02338\\Model: GPT5.6}
\date{19 July 2026}

\begin{document}
\maketitle

\begin{center}
\fbox{\parbox{0.92\textwidth}{\textbf{Status: candidate full solution, likely valid.}
The compact-operator-algebra subquestion on page~40 of Guo--Wang is answered
completely under the standard meaning of Banach-algebra isomorphism.  The
surrounding maximal Roe-algebra and $K$-theory questions are not claimed.}}
\end{center}

\section{Source and exact target}

Guo and Wang's paper \cite{GW} asks several questions about two maximal
$L^p$-Roe completions.  In the discussion following Question~8.1, on source
PDF page~40, the authors write that it is unknown to them whether
\[
  \Kop(L^p(\mu)) \quad\hbox{is isomorphic to}\quad \Kop(\ell^p),
\]
after noting the non-isomorphism of $L^p(0,1)$ and $\ell^p$.  We answer the
concrete case $L^p(0,1)$, interpreting ``isomorphic'' in the standard
Banach-algebra sense: a bounded linear bijective algebra homomorphism.

\begin{figure}[h]
\centering
\includegraphics[width=\textwidth]{figures/open_problem_crop.png}
\caption{The source paragraph on PDF page~40.  The relevant compact-operator
algebra sentence is the first complete sentence in the second paragraph.}
\end{figure}

\section{Proof intuition}

The algebra $\Kop(X)$ already contains a canonical copy of $X$, up to
isomorphism.  Choose a rank-one idempotent $P=x_0\otimes x_0^*$.  The left
ideal $\Kop(X)P$ consists exactly of the rank-one operators
$x\otimes x_0^*$ and is therefore isomorphic to $X$.  It is algebraically a
minimal nonzero left ideal.  An algebra isomorphism sends it to a minimal
left ideal $\Kop(Y)Q$.  Since $Q$ is a compact idempotent, it has finite rank;
minimality forces that rank to be one.  Thus the image ideal is isomorphic to
$Y$, recovering $X\cong Y$.

For the $L^p$ application, $L^p(0,1)\not\cong\ell^p$ when $p\ne2$.  At
$p=1$ the Schur property separates them.  For $1<p<\infty$, the Rademacher
functions span a complemented copy of $\ell^2$ in $L^p(0,1)$, while Pitt's
compactness theorem prevents the corresponding embedding or complementation
inside $\ell^p$ unless $p=2$.

\section{The algebraic rigidity theorem}

For Banach spaces $X,Y$, write $\Kop(X)$ for the Banach algebra of compact
operators on $X$ with the operator norm.

\begin{theorem}[compact-operator algebra rigidity]\label{thm:rigidity}
Let $X$ and $Y$ be nonzero Banach spaces over the same scalar field.  Then
\[
  \Kop(X)\cong\Kop(Y)\quad\text{as Banach algebras}
  \qquad\Longleftrightarrow\qquad
  X\cong Y\quad\text{as Banach spaces}.
\]
No approximation-property assumption is required.
\end{theorem}

\begin{proof}
Suppose first that $S:X\to Y$ is a Banach-space isomorphism.  Conjugation,
\[
  T\longmapsto STS^{-1},
\]
is a bounded Banach-algebra isomorphism from $\Kop(X)$ onto $\Kop(Y)$.

Conversely, let $\Phi:\Kop(X)\to\Kop(Y)$ be a bounded Banach-algebra
isomorphism.  Choose $x_0\in X$ and $x_0^*\in X^*$ with
$x_0^*(x_0)=1$, and put
\[
  P=x_0\otimes x_0^*,\qquad P(x)=x_0^*(x)x_0.
\]
Then $P$ is a rank-one idempotent.  Its generated left ideal is
\[
  \Kop(X)P=\{x\otimes x_0^*:x\in X\}.
\]
Indeed, $TP=(Tx_0)\otimes x_0^*$; conversely, every $x$ occurs as $Tx_0$
for a rank-one operator $T$.  The map
\[
  X\longrightarrow\Kop(X)P,\qquad x\longmapsto x\otimes x_0^*,
\]
is a Banach-space isomorphism, since
$\|x\otimes x_0^*\|=\|x\|\,\|x_0^*\|$.

The left ideal $\Kop(X)P$ is algebraically minimal.  To see this, take a
nonzero $x\otimes x_0^*$ in a left subideal.  Given any $z\in X$, a rank-one
operator $A$ can be chosen with $Ax=z$, so
$A(x\otimes x_0^*)=z\otimes x_0^*$ also lies in that subideal.

Set $Q=\Phi(P)$.  Then $Q$ is a nonzero compact idempotent and
\[
  \Phi(\Kop(X)P)=\Kop(Y)Q.
\]
Thus $\Kop(Y)Q$ is a minimal nonzero left ideal.  The range of a compact
idempotent is finite-dimensional: the idempotent restricts to the identity on
its range, and a compact identity can occur only in finite dimension.  If
$\operatorname{rank}Q>1$, choose a rank-one idempotent $R$ whose range lies
in $\operatorname{ran}Q$ and such that $RQ=R$.  Then
\[
  0\ne\Kop(Y)R\subsetneq\Kop(Y)Q,
\]
where properness follows because equality would express $Q=AR$ for some
$A$, forcing $\operatorname{rank}Q\le1$.  This contradicts minimality.
Therefore $Q$ has rank one.

Writing $Q=y_0\otimes y_0^*$ with $y_0^*(y_0)=1$, the same argument as above
identifies $\Kop(Y)Q$ with $Y$.  The bounded restriction of $\Phi$, with its
bounded inverse, now gives
\[
  X\cong\Kop(X)P\cong\Kop(Y)Q\cong Y.
\]
\end{proof}

\section{Application to \texorpdfstring{$L^p(0,1)$}{Lp(0,1)} and
\texorpdfstring{$\ell^p$}{ell-p}}

\begin{lemma}\label{lem:lp}
For $1\le p<\infty$,
\[
  L^p(0,1)\cong\ell^p \qquad\Longleftrightarrow\qquad p=2.
\]
\end{lemma}

\begin{proof}
For $p=2$, both spaces are separable infinite-dimensional Hilbert spaces and
are isometrically isomorphic.

For $p=1$, $\ell^1$ has the Schur property, whereas $L^1(0,1)$ does not.  In
fact, the Rademacher functions $r_n$ satisfy $\|r_n\|_1=1$ and converge weakly
to zero: for $g\in L^\infty(0,1)\subset L^2(0,1)$, the coefficients
$\int g r_n$ tend to zero by Bessel's inequality.

Let $1<p<\infty$.  Khintchine's inequality \cite{Khintchine} shows that the
closed Rademacher span
\[
  R_p=\overline{\operatorname{span}}\{r_n:n\ge1\}\subset L^p(0,1)
\]
is isomorphic to $\ell^2$.  It is complemented.  Indeed, for finite scalar
sequences $(b_n)$, H\"older's and Khintchine's inequalities give
\[
 \left|\sum_n b_n\int f r_n\right|
 \le \|f\|_p\left\|\sum_n b_nr_n\right\|_{p'}
 \le C_{p'}\|f\|_p\|(b_n)\|_2.
\]
Thus $(\int fr_n)_n\in\ell^2$ with norm controlled by $\|f\|_p$, and another
application of Khintchine's inequality makes
\[
  f\longmapsto\sum_{n=1}^\infty\left(\int_0^1 f r_n\right)r_n
\]
a bounded projection onto $R_p$.

Pitt's theorem \cite{Pitt} says that every bounded operator
$\ell^s\to\ell^r$ is compact when $1\le r<s<\infty$.  If $1<p<2$, an
embedding $\ell^2\hookrightarrow\ell^p$ would itself be compact by Pitt, which
is impossible for an isomorphic embedding of an infinite-dimensional Banach
space.  If $2<p<\infty$ and $\ell^p$ contained a complemented copy of
$\ell^2$, an embedding $J:\ell^2\to\ell^p$ and projection onto $J\ell^2$
would yield a bounded operator $S:\ell^p\to\ell^2$ with $SJ=I_{\ell^2}$.
Pitt makes $S$ compact, hence makes the identity compact, again impossible.
Therefore $\ell^p$ cannot be isomorphic to $L^p(0,1)$ for $p\ne2$.
\end{proof}

\begin{corollary}[answer to the source subquestion]\label{cor:answer}
For $1\le p<\infty$,
\[
  \boxed{\quad
  \Kop(L^p(0,1))\cong\Kop(\ell^p)\text{ as Banach algebras}
  \quad\Longleftrightarrow\quad p=2.
  \quad}
\]
In particular, the answer is negative for every $p\ne2$ and affirmative for
$p=2$.
\end{corollary}

\begin{proof}
Combine Theorem~\ref{thm:rigidity} with Lemma~\ref{lem:lp}.
\end{proof}

\section{Verification, limitations, and novelty check}

\paragraph{Verification.}
The proof was checked directly at each algebraic step.  The only external
inputs in Lemma~\ref{lem:lp} are Khintchine's inequality and Pitt's theorem;
the complemented Rademacher projection is derived explicitly.  The theorem
uses neither a basis nor an approximation property.  The rank argument also
checks that a minimal left ideal generated by a compact idempotent must come
from a rank-one idempotent.

\paragraph{Scope.}
The conclusion concerns bounded Banach-algebra isomorphisms.  It does not
decide whether the two compact-operator spaces are isomorphic merely as Banach
spaces.  It also does not decide whether the two maximal $L^p$-Roe
completions in the source coincide or have the same $K$-theory.  Finally, for
a general measure $\mu$, the answer depends on the Banach-space isomorphism
type of $L^p(\mu)$; the displayed corollary is for the nonatomic model
$L^p(0,1)$ explicitly named by the source.

\paragraph{Novelty search.}
On 19 July 2026, bounded searches covered the run's registry, solution,
attempt, and proof-gap indexes; exact arXiv-id/title searches; and web/arXiv
queries for variants of ``$\Kop(X)$ isomorphic to $\Kop(Y)$'', compact-operator
algebras and minimal left ideals, and
``$\Kop(L^p)$ isomorphic to $\Kop(\ell^p)$''.  No source stating this exact
answer was located.  The minimal-left-ideal mechanism is elementary and may
well be standard ring-theoretic folklore, so originality confidence is
conservative.  The mathematical claim does not depend on novelty.

\paragraph{Human review recommendation.}
Check that ``isomorphic'' in the source paragraph was intended as
Banach-algebra isomorphism rather than Banach-space isomorphism.  Under that
standard algebraic interpretation, the proof is short and the result should
be promoted after review.

\begin{thebibliography}{9}
\bibitem{GW}
L. Guo and Q. Wang,
\emph{Geometric Banach property (T) for metric spaces via Banach
representations of Roe algebras}, arXiv:2505.02338v2 (2025).

\bibitem{Khintchine}
A. Khintchine,
\emph{\"Uber dyadische Br\"uche},
Math. Z. \textbf{18} (1923), 109--116.

\bibitem{Pitt}
H. R. Pitt,
\emph{A note on bilinear forms},
J. London Math. Soc. \textbf{s1-11} (1936), 174--180.
\end{thebibliography}

\end{document}
